% ---------------------------------------------------------------------------
%  Section 3 -- Method
%
%  FIGURES EXPECTED (drop into figures/, names can be changed):
%    forward.pdf    -- forward diffusion in latent space   -> Fig.~\ref{fig:forward}
%    framework.pdf  -- full pipeline / reverse process     -> Fig.~\ref{fig:framework}
%
%  CITATION KEYS reused from the MedSegLatDiff bibliography:
%    ddpm, ddim, vae, ldm
%  CITATION KEYS supplied in refs_method.bib:
%    probunet, phiseg, mcdropout, deepensembles, cover2006elements
% ---------------------------------------------------------------------------

\section{Method}
\label{sec:method}

% ---------------------------------------------------------------------------
\subsection{Problem Formulation}
\label{subsec:problem}

We formulate brain tumour segmentation from incomplete multi-sequence MRI as
conditional generation in a learned latent space. Let
\begin{equation}
    \mathcal{D} = \bigl\{ (X^{(i)}, Y^{(i)}) \bigr\}_{i=1}^{N}
\end{equation}
denote a dataset of multi-sequence brain MRI studies, where
$X^{(i)} \in \mathbb{R}^{H \times W \times D \times M}$ is a co-registered
volume with $M = 4$ sequences (FLAIR, T1ce, T1, T2) and
$Y^{(i)} \in \{0,1,2,4\}^{H \times W \times D}$ is the voxel-wise expert
annotation. We target the \emph{whole tumour} (WT), the union of the necrotic
core, peritumoural oedema and enhancing tumour, giving the binary mask
\begin{equation}
    S = \mathbb{1}\bigl[\, Y > 0 \,\bigr] \in \{0,1\}^{H \times W \times D}.
    \label{eq:wt}
\end{equation}
A single binary region is the setting in which probabilistic segmentation is
conventionally studied and evaluated~\cite{probunet,phiseg},
and it keeps the uncertainty estimates of Section~\ref{subsec:uncertainty}
interpretable: dispersion over one foreground class has an unambiguous
meaning, whereas dispersion aggregated over overlapping regions of very
different size and difficulty does not.

\paragraph{Modality availability.}
In routine practice the four sequences are frequently incomplete: acquisitions
are shortened, corrupted by motion, or omitted by protocol. We model availability
explicitly with a binary vector $\mathbf{m} \in \{0,1\}^{M} \setminus
\{\mathbf{0}\}$, giving $2^{M} - 1 = 15$ distinct configurations, and define
the observed volume as
\begin{equation}
    X_{\mathbf{m}} = X \odot \mathbf{m},
    \label{eq:masking}
\end{equation}
where $\odot$ broadcasts $\mathbf{m}$ across the spatial dimensions so that
absent sequences are set to zero. Equation~\eqref{eq:masking} is a
deterministic map that can only destroy information, which is the property
Section~\ref{subsec:uncertainty} relies on.

\paragraph{Objective.}
Rather than learning a point estimate $f: X_{\mathbf{m}} \mapsto S$, we learn
the conditional distribution $p_{\theta}(S \mid X_{\mathbf{m}})$, from which
multiple plausible annotations can be drawn. A deterministic model returns one
boundary whatever the input contains, whereas a distribution can widen when
the evidence supporting that boundary is removed.

% ---------------------------------------------------------------------------
\subsection{Volumetric Perceptual Compression}
\label{subsec:compression}

Operating a diffusion model directly on $128^{3}$ volumes is prohibitive: the
denoising network must be evaluated once per sampling step, and drawing
several samples per case multiplies that cost again. We therefore follow the
latent diffusion paradigm~\cite{ldm} and move both the image and the
annotation into a compressed space before diffusion.

We train two separate variational autoencoders~\cite{vae}. Their inputs have
markedly different statistics. One is continuous and scanner-dependent, the
other is piecewise-constant binary structure, and a single shared encoder
would have to compromise between them. Each encoder
reduces the spatial resolution by a factor of four per axis, from $128^{3}$ to
$32^{3}$, a $64\times$ reduction in the number of spatial positions. In two dimensions this
compression mainly saves time; in three it is what keeps the model trainable
within a fixed memory budget.

\paragraph{Image autoencoder.}
The encoder $\mathcal{E}_{X}$ maps the masked volume $X_{\mathbf{m}}$ to a
posterior $\mathcal{N}\bigl(\mu_{X}, \sigma_{X}^{2}\bigr)$ over latents, and we
take its mean $z_{X} = \mu_{X} \in \mathbb{R}^{32^{3} \times C_{X}}$ as the
conditioning signal. Training minimises reconstruction error against the
observed sequences together with a Kullback--Leibler term,
\begin{equation}
    \mathcal{L}_{X}
    = \bigl\| X_{\mathbf{m}} - \mathcal{D}_{X}(z_{X}) \bigr\|_{2}^{2}
    + \beta_{X}\, D_{\mathrm{KL}}\!\left(
        \mathcal{N}(\mu_{X}, \sigma_{X}^{2}) \,\|\, \mathcal{N}(0, I)
      \right).
\end{equation}
Modality dropout is applied during autoencoder training as well as during
diffusion training, so that the encoder learns to represent partial inputs
rather than encountering them for the first time at inference.

\paragraph{Annotation autoencoder.}
The second autoencoder compresses the binary mask $S$ of Eq.~\eqref{eq:wt}.
Reconstruction combines a weighted binary cross-entropy with a soft Dice
term,
\begin{equation}
    \mathcal{L}_{S}
    = \underbrace{-\Bigl[\, w_{+}\, S \log \hat{S}
        + (1 - S) \log \bigl(1 - \hat{S}\bigr) \Bigr]}_{\text{weighted BCE}}
    \;+\; \underbrace{\left( 1 - \frac{2\,\mathrm{TP}}
        {2\,\mathrm{TP} + \mathrm{FP} + \mathrm{FN}} \right)}_{\text{soft Dice}}
    \;+\; \beta_{S}\, D_{\mathrm{KL}},
    \label{eq:mask-loss}
\end{equation}
where $\hat{S} = \mathcal{D}_{S}(z_{S})$, the weight $w_{+} > 1$ upweights
tumour voxels, and $\mathrm{TP}$, $\mathrm{FP}$ and $\mathrm{FN}$ are computed
from $\hat{S}$ and $S$ without thresholding. Both terms address the same
problem from different directions. Tumour occupies a small fraction of each
volume, so an unweighted per-voxel objective drives the decoder towards the
empty mask, which has low average error but no practical value; the weighting
penalises missed tumour voxels more heavily, while the Dice term optimises
region overlap directly rather than voxel-wise agreement. The resulting latent
$z_{S} \in \mathbb{R}^{32^{3} \times C_{S}}$ is the variable the diffusion
model generates.

Both autoencoders are frozen after this stage.

% ---------------------------------------------------------------------------
\subsection{Latent Diffusion with Modality Dropout}
\label{subsec:diffusion}

\begin{figure}[t]
    \centering
    \includegraphics[width=\textwidth]{forward.pdf}
    \caption{Forward diffusion in latent space. Gaussian noise is added to the
    encoded annotation $z_{S,0}$ over $T$ steps until the representation is
    indistinguishable from $\mathcal{N}(0, I)$. The encoded image is not
    noised; it conditions the reverse process.}
    \label{fig:forward}
\end{figure}

\paragraph{Forward process.}
Diffusion is applied to the annotation latent $z_{S,0} = \mathcal{E}_{S}(S)$.
Following DDPM~\cite{ddpm}, noise is added according to
\begin{equation}
    q\bigl(z_{S,t} \mid z_{S,t-1}\bigr)
    = \mathcal{N}\bigl(z_{S,t};\, \sqrt{\alpha_{t}}\, z_{S,t-1},\,
                       (1-\alpha_{t}) I \bigr),
\end{equation}
which admits the closed form
\begin{equation}
    q\bigl(z_{S,t} \mid z_{S,0}\bigr)
    = \mathcal{N}\bigl(z_{S,t};\, \sqrt{\gamma_{t}}\, z_{S,0},\,
                       (1-\gamma_{t}) I \bigr),
    \qquad \gamma_{t} = \prod_{i=1}^{t} \alpha_{i},
\end{equation}
allowing any timestep to be sampled directly during training. The process is
illustrated in Fig.~\ref{fig:forward}.

\paragraph{Conditioning.}
The denoising network is conditioned on the encoded image by channel-wise
concatenation in the latent space,
\begin{equation}
    z_{S,t}^{\mathrm{cond}} = z_{S,t} \oplus \pi\bigl(z_{X}\bigr),
\end{equation}
where $\pi$ is a learned $1\times1\times1$ projection that reduces the image
latent to a smaller channel count before concatenation, keeping the
conditioning from dominating the input by sheer width. Because the two latents
share a spatial resolution of $32^{3}$, conditioning requires no resampling or
cross-attention.

\paragraph{Modality dropout.}
The availability vector $\mathbf{m}$ is resampled independently for every
training example, uniformly over the $15$ non-empty configurations. A single
set of parameters therefore covers every acquisition protocol, in contrast to
training a separate model per configuration or imputing absent sequences
before segmentation. Uniform sampling deliberately over-represents degraded
inputs relative to their clinical frequency, so that the sparse configurations,
where uncertainty matters most, are not treated as rare corner cases.

\paragraph{Training objective.}
With $\epsilon \sim \mathcal{N}(0, I)$ and $t$ drawn uniformly over
$\{1, \dots, T\}$, the network $\epsilon_{\theta}$ is trained to predict the
noise added at step $t$:
\begin{equation}
    \mathcal{L}_{\mathrm{diff}}
    = \mathbb{E}_{z_{S,0},\, \mathbf{m},\, \epsilon,\, t}
      \Bigl[\, \bigl\| \epsilon
        - \epsilon_{\theta}\bigl(z_{S,t}^{\mathrm{cond}}, t\bigr)
      \bigr\|_{2}^{2} \Bigr].
    \label{eq:diffusion-loss}
\end{equation}
The expectation over $\mathbf{m}$ in Eq.~\eqref{eq:diffusion-loss} is what
distinguishes this objective from a standard conditional diffusion loss: the
model is optimised over the whole family of degraded inputs rather than a
single fixed one.

% ---------------------------------------------------------------------------
\subsection{Sampling and Predictive Uncertainty}
\label{subsec:uncertainty}

\begin{figure}[t]
    \centering
    \includegraphics[width=\textwidth]{framework.pdf}
    \caption{Inference. The available sequences are encoded once; $N$
    independent noise draws are denoised in parallel and decoded, yielding $N$
    plausible annotations. Their mean is the prediction and their per-voxel
    variance is the uncertainty estimate.}
    \label{fig:framework}
\end{figure}

At inference, the available sequences are encoded once and $N$ independent
initial noise tensors $z_{S,T}^{1}, \dots, z_{S,T}^{N} \sim \mathcal{N}(0, I)$
are denoised with DDIM~\cite{ddim} under the same conditioning. Decoding each
trajectory gives $N$ annotations $S^{i} = \mathcal{D}_{S}(z_{S,0}^{i})$ for
$i = 1, \dots, N$, which are samples from the learned posterior
$p_{\theta}(S \mid X_{\mathbf{m}})$. Their per-voxel mean and variance,
\begin{equation}
    \mu(v) = \frac{1}{N} \sum_{i=1}^{N} S^{i}(v),
    \qquad
    \sigma^{2}(v) = \frac{1}{N} \sum_{i=1}^{N}
        \bigl( S^{i}(v) - \mu(v) \bigr)^{2},
    \label{eq:mean-var}
\end{equation}
give the prediction and the uncertainty map respectively; the final
segmentation is $\mathbb{1}[\mu(v) > \tau]$ with $\tau = 0.5$.

The dispersion in Eq.~\eqref{eq:mean-var} is not noise added to produce
variability. The model is trained to sample from
$p(S \mid X_{\mathbf{m}})$, so the spread across draws is the posterior the
model has learned. This differs from Monte Carlo dropout~\cite{mcdropout},
which perturbs the network itself and reports the resulting instability, and
from deep ensembles~\cite{deepensembles}, which capture parameter uncertainty
but require training several models.

\paragraph{Restricting the average.}
Summarising an uncertainty map by one number requires care. The mean of
$\sigma^{2}$ over the whole volume is dominated by background: the great
majority of voxels are unambiguously healthy tissue, every sample agrees, and
their variance is identically zero. Such an average therefore scales inversely
with lesion size rather than with confidence, and it is not comparable across
modality configurations, because the predicted lesion extent itself depends
on $\mathbf{m}$. We instead average only over the
region in which the samples are able to disagree: the set $\Omega$ of voxels
that are annotated tumour, or that at least one sample labels tumour. Writing
$|\Omega|$ for its size,
\begin{equation}
    U\bigl(X_{\mathbf{m}}\bigr)
    = \frac{1}{|\Omega|} \sum_{v \in \Omega} \sigma^{2}(v).
    \label{eq:uncertainty}
\end{equation}
Equation~\eqref{eq:uncertainty} measures how much plausible annotations
disagree in the region where they can disagree, and is the only form of the
quantity that stays comparable across the configurations we compare below.

\paragraph{Expected behaviour under information loss.}
The masking in Eq.~\eqref{eq:masking} is a deterministic function of $X$, so
the data processing inequality~\cite{cover2006elements} gives, for any nested
pair of configurations $\mathbf{m}' \subseteq \mathbf{m}$,
\begin{equation}
    I\bigl(S; X_{\mathbf{m}'}\bigr) \;\le\; I\bigl(S; X_{\mathbf{m}}\bigr),
    \qquad\text{hence}\qquad
    H\bigl(S \mid X_{\mathbf{m}'}\bigr)
        \;\ge\; H\bigl(S \mid X_{\mathbf{m}}\bigr).
    \label{eq:dpi}
\end{equation}
Removing a sequence cannot create information about the annotation, so the
true posterior is at least as dispersed after removal as before, and a
well-calibrated model should report uncertainty that follows the same
ordering. Two qualifications apply. The inequality holds only for
\emph{nested} configurations, so whether T1CE is more informative than T1 is
an empirical question rather than a consequence of the bound. It also
constrains the true posterior and not the model's estimate of it, so a poorly
calibrated model may stay confident while the supporting evidence disappears.

Since every case is evaluated under all $15$ configurations, this effect can
be measured within each case. For case $i$ we report
\begin{equation}
    \Delta U_{i}(\mathbf{m})
    = U\bigl(X^{(i)}_{\mathbf{m}}\bigr) - U\bigl(X^{(i)}_{\mathbf{1}}\bigr),
    \label{eq:delta-u}
\end{equation}
where $\mathbf{1}$ denotes the fully observed input. Lesion size and case
difficulty are shared between the two terms and cancel, leaving the
contribution of the removed sequences.